\documentclass[12pt,a4paper,oneside]{article}

\usepackage[utf8]{inputenc}
\usepackage[T1]{fontenc}
\usepackage[brazil]{babel}
\usepackage{geometry}
\geometry{a4paper, margin=2.5cm}
\usepackage{indentfirst}
\usepackage{times}
\usepackage{graphicx}
\usepackage{hyperref}
\usepackage{abntex2cite}

\setlength{\parindent}{1.25cm}
\setlength{\parskip}{0.3cm}

\begin{document}

\begin{center}
  \textbf{\large ESCOLA E FACULDADE DE TECNOLOGIA SENAI GASPAR RICARDO JÚNIOR}\\
  \textbf{Desenvolvimento de Sistemas}\\[1cm]
  \textbf{\Large Lucas Miguel Leite}\\[0.5cm]
  \textbf{Professor: Deivison Takatu}\\[2cm]
  \textbf{\LARGE\textbf{O Papel da TI nas Organizações Brasileiras: Análise da 36ª Pesquisa FGVcia de 2025}}\\[1cm]
  Sorocaba -- 18/03/2026
\end{center}

\vspace{1cm}

\begin{center}
  \textbf{RESUMO}
\end{center}

Este artigo analisa o papel estratégico da Tecnologia da Informação nas organizações brasileiras com base nos indicadores da 36ª Pesquisa FGVcia de 2025, que analisou 2.672 empresas abordando infraestrutura, investimentos e sistemas. Três indicadores são examinados em profundidade: o alto nível de investimento em TI (10\% da receita das empresas), a centralidade da Inteligência Artificial (destaque para Microsoft Copilot, ChatGPT e Google Gemini) e a consolidação do processamento em nuvem (52\% do processamento). A análise relaciona esses dados aos conceitos de alinhamento estratégico, transformação digital e vantagem competitiva, concluindo que a TI é hoje um motor indispensável para o modelo de negócios atual e que a sobrevivência corporativa depende da integração eficiente de sistemas, dados e automação inteligente.

\vspace{0.5cm}
\textbf{Palavras-chave:} TI Estratégica. FGVcia. Inteligência Artificial. Computação em Nuvem. Transformação Digital.

\newpage

\section{Introdução}

A 36ª Pesquisa FGVcia de 2025, com histórico de 36 anos e análise de 2.672 empresas, constitui uma das mais abrangentes fontes de dados sobre o uso da tecnologia nas organizações brasileiras. A pesquisa aborda infraestrutura, investimentos e sistemas, oferecendo indicadores que permitem avaliar o papel estratégico da TI no cenário nacional \cite{fgvcia2025}.

\section{Indicador 1: Alto Nível de Investimento em TI}

O gasto e investimento em TI atingiu a marca de 10\% da receita das empresas, número que continua crescendo. O custo anual médio de TI por usuário nas empresas pesquisadas foi de R\$ 60.000. Esse nível de investimento demonstra que a tecnologia deixou de ser uma área de suporte puramente operacional, evidenciando um foco no Alinhamento Estratégico, onde a TI é fundamental para o atingimento das metas organizacionais \cite{rocha2018}.

\section{Indicador 2: Centralidade da Inteligência Artificial}

A Inteligência Artificial aparece como um dos principais projetos de TI nas empresas. Na categoria inédita de IA Generativa, destacam-se o Microsoft Copilot (40\%), o ChatGPT (32\%) e o Google Gemini (20\%), sendo utilizados para Chatbot, Machine Learning e Reconhecimento Biométrico. A adoção acelerada reflete o esforço das organizações em impulsionar a Transformação Digital, gerando soluções que integram e ampliam a capacidade humana com a digital \cite{kluyver2007}.

\section{Indicador 3: Consolidação do Processamento em Nuvem}

A Nuvem já responde, em média, por 52\% do processamento nas empresas pesquisadas. Esse número indica uma busca por maior flexibilidade, redução de custos com infraestrutura física e escalabilidade. Especialmente nas maiores empresas, a Nuvem aparece como um dos projetos principais, apoiando inclusive pautas de ESG \cite{fgvcia2025}.

\section{Análise e Relação com os Conceitos Estratégicos}

A partir dos indicadores selecionados, é possível estabelecer uma forte relação com o papel estratégico e transformador da TI nas organizações contemporâneas.

\subsection{A TI como Pilar Estratégico}

Os gastos e investimentos em TI estão crescendo em valor, maturidade e importância nos negócios. Isso evidencia que as empresas brasileiras reconhecem a TI como instrumento essencial de competitividade e não mais como centro de custo.

\subsection{Adoção de IA e Vantagem Competitiva}

Os vários tipos de Inteligência Artificial permanecem no centro das pautas empresariais, acompanhados por Segurança e Inteligência Analítica. A adoção de IA Generativa gera soluções que mantêm a competitividade em mercados cada vez mais dinâmicos.

\subsection{Infraestrutura Flexível e Escalável}

A transição de mais da metade do processamento para a Nuvem indica maturidade na adoção de infraestrutura moderna, proporcionando agilidade e resiliência operacional.

\section{Conclusão}

Os dados da FGVcia evidenciam que a TI é hoje um motor indispensável para o modelo de negócios atual brasileiro. O foco permanente está no Alinhamento Estratégico e na Transformação Digital, provando que a sobrevivência corporativa depende da integração eficiente de sistemas, dados (Nuvem) e automação inteligente (IA). Os indicadores de investimento elevado, centralidade da IA e consolidação da nuvem convergem para um cenário em que a TI deixou definitivamente de ser suporte para se tornar o núcleo da estratégia organizacional \cite{porter1989}.

\bibliographystyle{plain}
\bibliography{referencias}

\end{document}
